\documentclass[11pt]{article}
\usepackage[margin=1in]{geometry}
\usepackage{amsmath,amssymb,microtype}
\usepackage[colorlinks=true,urlcolor=blue,citecolor=blue]{hyperref}
\title{Solel's TRO-range conjecture was refuted by arXiv:math/0510641}
\author{Literature-status note for \texttt{fa\_banach\_001}}
\date{26 August 2026}

\begin{document}
\maketitle

\noindent\textbf{Status.}
\texttt{literature\_already\_answered\_negative}.  No new mathematical
result is claimed by this packet.

\section*{Original conjecture}

Let $\mathcal D\subseteq B(\mathcal H)$ be a maximal abelian selfadjoint
algebra (MASA).  Katavolos and Paulsen record on PDF page 2 of
arXiv:math/0301298 the conjecture of Solel that if
\[
 P:B(\mathcal H)\longrightarrow B(\mathcal H)
\]
is a contractive idempotent $\mathcal D$-bimodule map, then its range
$\mathcal M=P(B(\mathcal H))$ is a ternary ring of operators in the inherited
product:
\[
 \mathcal M\mathcal M^*\mathcal M\subseteq\mathcal M.
\]
Solel had proved this under the additional assumption that $P$ is
weak-star continuous.  The conjecture asked whether the same conclusion
holds without that normality assumption \cite{KP2003}.

\section*{Explicit later counterexample}

Paulsen's arXiv:math/0510641 says in its abstract and introduction that it
constructs a counterexample to precisely this conjecture \cite{P2005}.
The construction has two stages.

First, Theorem 2.15 on PDF page 10 produces a unital, completely positive,
idempotent, $\mathbb Z$-equivariant map
\[
 \Phi:\ell^\infty(\mathbb Z)\longrightarrow\ell^\infty(\mathbb Z)
\]
whose range is not a $C^*$-subalgebra.  The proof uses the
$\mathbb Z$-injective envelope of $C(\mathbb T)$ to obtain two distinct
equivariant embeddings and averages them; the resulting range fails
multiplicative closure.

Second, Theorem 3.1 on PDF page 12 lifts any such $\Phi$ to a unital
completely positive $L^\infty(\mathbb T)$-bimodule map
\[
 \Gamma:B(L^2(\mathbb T))\longrightarrow B(L^2(\mathbb T)),
\]
and preserves idempotence.  Corollary 3.2 on PDF pages 12--13 states that the
range of the lifted $\Gamma$ is not a $C^*$-subalgebra.

This refutes the conjecture exactly.  Complete positivity and unitality imply
$\|\Gamma\|=1$, so $\Gamma$ is contractive.  Its idempotent range contains
$I$.  If that range were a TRO, then for range elements $A,B$ the ternary
products
\[
 A I^* B=AB,\qquad I A^* I=A^*
\]
would make it a $C^*$-subalgebra, contrary to Corollary 3.2.  The construction
is non-weak-star-continuous, so Solel's positive normal-case theorem remains
intact.

\section*{Scope and search}

The four run indexes, both exact titles, and exact MASA/TRO conjecture phrases
were checked through 26 August 2026.  arXiv:math/0510641 is a direct explicit
answer, not an inferred consequence.  This packet does not claim to settle
the other questions in arXiv:math/0301298 concerning norm gaps, a general
symbol calculus, or the commutator ideals of the restriction maps.

\begin{thebibliography}{9}
\bibitem{KP2003}
A.~Katavolos and V.~I.~Paulsen,
\emph{On the ranges of bimodule projections},
Canad. Math. Bull. \textbf{48} (2005), 97--111;
arXiv:math/0301298,
\url{https://arxiv.org/abs/math/0301298}.

\bibitem{P2005}
V.~I.~Paulsen,
\emph{Equivariant Maps and Bimodule Projections},
J. Funct. Anal. \textbf{240} (2006), 495--507;
arXiv:math/0510641,
\url{https://arxiv.org/abs/math/0510641}.
\end{thebibliography}

\end{document}
